% ================= DANH MỤC TỪ VIẾT TẮT =================
\newpage
\thispagestyle{plain}

\begin{center}
    {\bfseries\Large DANH MỤC TỪ VIẾT TẮT}
\end{center}

\vspace{0.8cm}

\renewcommand{\arraystretch}{1.25}

% Quy ước: xếp theo THỨ TỰ BẢNG CHỮ CÁI.
% Dạng: Từ viết tắt & Dạng đầy đủ tiếng Anh -- Giải nghĩa tiếng Việt \\
\begin{longtable}{>{\centering\arraybackslash}p{3.2cm} p{10.5cm}}
    \toprule
    \textbf{Từ viết tắt} & \textbf{Diễn giải} \\
    \midrule
    \endfirsthead

    \toprule
    \textbf{Từ viết tắt} & \textbf{Diễn giải} \\
    \midrule
    \endhead

    AABB & Axis-Aligned Bounding Box -- Hộp bao song song trục toạ độ \\

    AP & Average Precision -- Độ chính xác trung bình \\

    CLIP & Contrastive Language--Image Pre-training -- Mô hình học tương phản ảnh--văn bản \\

    COCO & Common Objects in Context -- Bộ dữ liệu ảnh đời thường có chú thích \\

    CSV & Comma-Separated Values -- Định dạng dữ liệu phân tách bằng dấu phẩy \\

    GT & Ground Truth -- Nhãn thật dùng để đối chiếu \\

    IoU & Intersection over Union -- Tỉ số giao trên hợp \\

    LOO & Leave-One-Out -- Thủ tục để-lại-một khi đánh giá \\

    MAD & Median Absolute Deviation -- Độ lệch tuyệt đối trung vị \\

    MAE & Mean Absolute Error -- Sai số tuyệt đối trung bình \\

    mAP & mean Average Precision -- Độ chính xác trung bình bình quân \\

    MiDaS & Mixed Data Sampling -- Mô hình ước lượng độ sâu đơn mắt \\

    NMS & Non-Maximum Suppression -- Loại bỏ khung trùng lặp \\

    OBB & Oriented Bounding Box -- Hộp bao có định hướng \\

    OCID & Object Clutter Indoor Dataset -- Bộ dữ liệu RGB-D cảnh trong nhà lộn xộn \\

    RANSAC & Random Sample Consensus -- Thuật toán ước lượng mô hình bền với ngoại lai \\

    RGB-D & Red-Green-Blue + Depth -- Ảnh màu kèm ảnh độ sâu đã căn chỉnh \\

    RMSE & Root Mean Square Error -- Căn bậc hai của sai số bình phương trung bình \\

    SHA-256 & Secure Hash Algorithm 256-bit -- Hàm băm dùng để xác thực tệp trọng số \\

    YOLO & You Only Look Once -- Họ mô hình phát hiện đối tượng một giai đoạn \\

    \bottomrule
\end{longtable}
